\chapter{Drugs}

\section{Antibiotics}

\subsection{Penicillin}

\subsection{Amoxicillin}


\section{Proton Pump Inhibitors (PPIs)}

Examples include

\begin{itemize}
    \item Omeprazole
    \item Lansoprazole
    \item Esomeprazole
    \item Pantoprazole
\end{itemize}

PPIs travel through the stomach, and are absorbed in the small intestine into the blood. They act on gastric parietal cells to inhibit the release of gastric acid.

\textbf{Side effects} include:

\begin{itemize}
    \item Increased risk of \hyperref[condition:c-diff]{\textit{c. diff}} infection
\end{itemize}

\section{Alginates}

Examples include

\begin{itemize}
    \item Gaviscon
    \item Peptac
    \item Wind-eze
\end{itemize}

Forms a 'raft' over the stomach contents.

\section{Antacids}

\textbf{Mechanism of action}: neutralizes stomach acid with a chemical reaction.

\textbf{Examples include}:

\begin{itemize}
    \item Calcium carbonate (\ce{CaCO3})
    \item Magnesium carbonate (\ce{MgCO3})
\end{itemize}

For calcium carbonate, the reaction is as follows:

\begin{equation}
    \ce{CaCO3 + 2HCl -> CaCl2 + H2O + CO2}
\end{equation}

\section{Prokinetics}

\textbf{Examples include}:

\begin{itemize}
    \item Metoclopramide
    \item Domperidone
    \item Enythromycin (off-license)
\end{itemize}

\section{H2-receptor antagonists (H2RAs)}

\textbf{Mechanism of action}: binds to H2 receptors on gastric parietal cells, preventing histamine from binding and thus reducing gastric acid secretion.

\textbf{Examples include}:

\begin{itemize}
    \item Famotidine
    \item Cimetidine
\end{itemize}

\section{Non-Steroidal Anti-Inflammatory Drugs (NSAIDs)}

\textbf{Examples include}:

\begin{itemize}
    \item Ibuprofen
    \item Naproxen
    \item Diclofenac
\end{itemize}

\textbf{Mechanism of action}

Prostaglandins produce pain/inflammation/fever. If we inhibit COX-2, we can avoid producing prostaglandins, but as a side effect that prevents the protection of the gastric mucosa.